% Options for packages loaded elsewhere
% Options for packages loaded elsewhere
\PassOptionsToPackage{unicode}{hyperref}
\PassOptionsToPackage{hyphens}{url}
\PassOptionsToPackage{dvipsnames,svgnames,x11names}{xcolor}
%
\documentclass[
  letterpaper,
  DIV=11,
  numbers=noendperiod]{scrartcl}
\usepackage{xcolor}
\usepackage{amsmath,amssymb}
\setcounter{secnumdepth}{-\maxdimen} % remove section numbering
\usepackage{iftex}
\ifPDFTeX
  \usepackage[T1]{fontenc}
  \usepackage[utf8]{inputenc}
  \usepackage{textcomp} % provide euro and other symbols
\else % if luatex or xetex
  \usepackage{unicode-math} % this also loads fontspec
  \defaultfontfeatures{Scale=MatchLowercase}
  \defaultfontfeatures[\rmfamily]{Ligatures=TeX,Scale=1}
\fi
\usepackage{lmodern}
\ifPDFTeX\else
  % xetex/luatex font selection
\fi
% Use upquote if available, for straight quotes in verbatim environments
\IfFileExists{upquote.sty}{\usepackage{upquote}}{}
\IfFileExists{microtype.sty}{% use microtype if available
  \usepackage[]{microtype}
  \UseMicrotypeSet[protrusion]{basicmath} % disable protrusion for tt fonts
}{}
\makeatletter
\@ifundefined{KOMAClassName}{% if non-KOMA class
  \IfFileExists{parskip.sty}{%
    \usepackage{parskip}
  }{% else
    \setlength{\parindent}{0pt}
    \setlength{\parskip}{6pt plus 2pt minus 1pt}}
}{% if KOMA class
  \KOMAoptions{parskip=half}}
\makeatother
% Make \paragraph and \subparagraph free-standing
\makeatletter
\ifx\paragraph\undefined\else
  \let\oldparagraph\paragraph
  \renewcommand{\paragraph}{
    \@ifstar
      \xxxParagraphStar
      \xxxParagraphNoStar
  }
  \newcommand{\xxxParagraphStar}[1]{\oldparagraph*{#1}\mbox{}}
  \newcommand{\xxxParagraphNoStar}[1]{\oldparagraph{#1}\mbox{}}
\fi
\ifx\subparagraph\undefined\else
  \let\oldsubparagraph\subparagraph
  \renewcommand{\subparagraph}{
    \@ifstar
      \xxxSubParagraphStar
      \xxxSubParagraphNoStar
  }
  \newcommand{\xxxSubParagraphStar}[1]{\oldsubparagraph*{#1}\mbox{}}
  \newcommand{\xxxSubParagraphNoStar}[1]{\oldsubparagraph{#1}\mbox{}}
\fi
\makeatother

\usepackage{color}
\usepackage{fancyvrb}
\newcommand{\VerbBar}{|}
\newcommand{\VERB}{\Verb[commandchars=\\\{\}]}
\DefineVerbatimEnvironment{Highlighting}{Verbatim}{commandchars=\\\{\}}
% Add ',fontsize=\small' for more characters per line
\usepackage{framed}
\definecolor{shadecolor}{RGB}{241,243,245}
\newenvironment{Shaded}{\begin{snugshade}}{\end{snugshade}}
\newcommand{\AlertTok}[1]{\textcolor[rgb]{0.68,0.00,0.00}{#1}}
\newcommand{\AnnotationTok}[1]{\textcolor[rgb]{0.37,0.37,0.37}{#1}}
\newcommand{\AttributeTok}[1]{\textcolor[rgb]{0.40,0.45,0.13}{#1}}
\newcommand{\BaseNTok}[1]{\textcolor[rgb]{0.68,0.00,0.00}{#1}}
\newcommand{\BuiltInTok}[1]{\textcolor[rgb]{0.00,0.23,0.31}{#1}}
\newcommand{\CharTok}[1]{\textcolor[rgb]{0.13,0.47,0.30}{#1}}
\newcommand{\CommentTok}[1]{\textcolor[rgb]{0.37,0.37,0.37}{#1}}
\newcommand{\CommentVarTok}[1]{\textcolor[rgb]{0.37,0.37,0.37}{\textit{#1}}}
\newcommand{\ConstantTok}[1]{\textcolor[rgb]{0.56,0.35,0.01}{#1}}
\newcommand{\ControlFlowTok}[1]{\textcolor[rgb]{0.00,0.23,0.31}{\textbf{#1}}}
\newcommand{\DataTypeTok}[1]{\textcolor[rgb]{0.68,0.00,0.00}{#1}}
\newcommand{\DecValTok}[1]{\textcolor[rgb]{0.68,0.00,0.00}{#1}}
\newcommand{\DocumentationTok}[1]{\textcolor[rgb]{0.37,0.37,0.37}{\textit{#1}}}
\newcommand{\ErrorTok}[1]{\textcolor[rgb]{0.68,0.00,0.00}{#1}}
\newcommand{\ExtensionTok}[1]{\textcolor[rgb]{0.00,0.23,0.31}{#1}}
\newcommand{\FloatTok}[1]{\textcolor[rgb]{0.68,0.00,0.00}{#1}}
\newcommand{\FunctionTok}[1]{\textcolor[rgb]{0.28,0.35,0.67}{#1}}
\newcommand{\ImportTok}[1]{\textcolor[rgb]{0.00,0.46,0.62}{#1}}
\newcommand{\InformationTok}[1]{\textcolor[rgb]{0.37,0.37,0.37}{#1}}
\newcommand{\KeywordTok}[1]{\textcolor[rgb]{0.00,0.23,0.31}{\textbf{#1}}}
\newcommand{\NormalTok}[1]{\textcolor[rgb]{0.00,0.23,0.31}{#1}}
\newcommand{\OperatorTok}[1]{\textcolor[rgb]{0.37,0.37,0.37}{#1}}
\newcommand{\OtherTok}[1]{\textcolor[rgb]{0.00,0.23,0.31}{#1}}
\newcommand{\PreprocessorTok}[1]{\textcolor[rgb]{0.68,0.00,0.00}{#1}}
\newcommand{\RegionMarkerTok}[1]{\textcolor[rgb]{0.00,0.23,0.31}{#1}}
\newcommand{\SpecialCharTok}[1]{\textcolor[rgb]{0.37,0.37,0.37}{#1}}
\newcommand{\SpecialStringTok}[1]{\textcolor[rgb]{0.13,0.47,0.30}{#1}}
\newcommand{\StringTok}[1]{\textcolor[rgb]{0.13,0.47,0.30}{#1}}
\newcommand{\VariableTok}[1]{\textcolor[rgb]{0.07,0.07,0.07}{#1}}
\newcommand{\VerbatimStringTok}[1]{\textcolor[rgb]{0.13,0.47,0.30}{#1}}
\newcommand{\WarningTok}[1]{\textcolor[rgb]{0.37,0.37,0.37}{\textit{#1}}}

\usepackage{longtable,booktabs,array}
\usepackage{calc} % for calculating minipage widths
% Correct order of tables after \paragraph or \subparagraph
\usepackage{etoolbox}
\makeatletter
\patchcmd\longtable{\par}{\if@noskipsec\mbox{}\fi\par}{}{}
\makeatother
% Allow footnotes in longtable head/foot
\IfFileExists{footnotehyper.sty}{\usepackage{footnotehyper}}{\usepackage{footnote}}
\makesavenoteenv{longtable}
\usepackage{graphicx}
\makeatletter
\newsavebox\pandoc@box
\newcommand*\pandocbounded[1]{% scales image to fit in text height/width
  \sbox\pandoc@box{#1}%
  \Gscale@div\@tempa{\textheight}{\dimexpr\ht\pandoc@box+\dp\pandoc@box\relax}%
  \Gscale@div\@tempb{\linewidth}{\wd\pandoc@box}%
  \ifdim\@tempb\p@<\@tempa\p@\let\@tempa\@tempb\fi% select the smaller of both
  \ifdim\@tempa\p@<\p@\scalebox{\@tempa}{\usebox\pandoc@box}%
  \else\usebox{\pandoc@box}%
  \fi%
}
% Set default figure placement to htbp
\def\fps@figure{htbp}
\makeatother


% definitions for citeproc citations
\NewDocumentCommand\citeproctext{}{}
\NewDocumentCommand\citeproc{mm}{%
  \begingroup\def\citeproctext{#2}\cite{#1}\endgroup}
\makeatletter
 % allow citations to break across lines
 \let\@cite@ofmt\@firstofone
 % avoid brackets around text for \cite:
 \def\@biblabel#1{}
 \def\@cite#1#2{{#1\if@tempswa , #2\fi}}
\makeatother
\newlength{\cslhangindent}
\setlength{\cslhangindent}{1.5em}
\newlength{\csllabelwidth}
\setlength{\csllabelwidth}{3em}
\newenvironment{CSLReferences}[2] % #1 hanging-indent, #2 entry-spacing
 {\begin{list}{}{%
  \setlength{\itemindent}{0pt}
  \setlength{\leftmargin}{0pt}
  \setlength{\parsep}{0pt}
  % turn on hanging indent if param 1 is 1
  \ifodd #1
   \setlength{\leftmargin}{\cslhangindent}
   \setlength{\itemindent}{-1\cslhangindent}
  \fi
  % set entry spacing
  \setlength{\itemsep}{#2\baselineskip}}}
 {\end{list}}
\usepackage{calc}
\newcommand{\CSLBlock}[1]{\hfill\break\parbox[t]{\linewidth}{\strut\ignorespaces#1\strut}}
\newcommand{\CSLLeftMargin}[1]{\parbox[t]{\csllabelwidth}{\strut#1\strut}}
\newcommand{\CSLRightInline}[1]{\parbox[t]{\linewidth - \csllabelwidth}{\strut#1\strut}}
\newcommand{\CSLIndent}[1]{\hspace{\cslhangindent}#1}



\setlength{\emergencystretch}{3em} % prevent overfull lines

\providecommand{\tightlist}{%
  \setlength{\itemsep}{0pt}\setlength{\parskip}{0pt}}



 


\KOMAoption{captions}{tableheading}
\usepackage{amsmath}
\usepackage{amssymb}
\makeatletter
\@ifpackageloaded{tcolorbox}{}{\usepackage[skins,breakable]{tcolorbox}}
\@ifpackageloaded{fontawesome5}{}{\usepackage{fontawesome5}}
\definecolor{quarto-callout-color}{HTML}{909090}
\definecolor{quarto-callout-note-color}{HTML}{0758E5}
\definecolor{quarto-callout-important-color}{HTML}{CC1914}
\definecolor{quarto-callout-warning-color}{HTML}{EB9113}
\definecolor{quarto-callout-tip-color}{HTML}{00A047}
\definecolor{quarto-callout-caution-color}{HTML}{FC5300}
\definecolor{quarto-callout-color-frame}{HTML}{acacac}
\definecolor{quarto-callout-note-color-frame}{HTML}{4582ec}
\definecolor{quarto-callout-important-color-frame}{HTML}{d9534f}
\definecolor{quarto-callout-warning-color-frame}{HTML}{f0ad4e}
\definecolor{quarto-callout-tip-color-frame}{HTML}{02b875}
\definecolor{quarto-callout-caution-color-frame}{HTML}{fd7e14}
\makeatother
\makeatletter
\@ifpackageloaded{caption}{}{\usepackage{caption}}
\AtBeginDocument{%
\ifdefined\contentsname
  \renewcommand*\contentsname{Table of contents}
\else
  \newcommand\contentsname{Table of contents}
\fi
\ifdefined\listfigurename
  \renewcommand*\listfigurename{List of Figures}
\else
  \newcommand\listfigurename{List of Figures}
\fi
\ifdefined\listtablename
  \renewcommand*\listtablename{List of Tables}
\else
  \newcommand\listtablename{List of Tables}
\fi
\ifdefined\figurename
  \renewcommand*\figurename{Figure}
\else
  \newcommand\figurename{Figure}
\fi
\ifdefined\tablename
  \renewcommand*\tablename{Table}
\else
  \newcommand\tablename{Table}
\fi
}
\@ifpackageloaded{float}{}{\usepackage{float}}
\floatstyle{ruled}
\@ifundefined{c@chapter}{\newfloat{codelisting}{h}{lop}}{\newfloat{codelisting}{h}{lop}[chapter]}
\floatname{codelisting}{Listing}
\newcommand*\listoflistings{\listof{codelisting}{List of Listings}}
\makeatother
\makeatletter
\makeatother
\makeatletter
\@ifpackageloaded{caption}{}{\usepackage{caption}}
\@ifpackageloaded{subcaption}{}{\usepackage{subcaption}}
\makeatother
\usepackage{bookmark}
\IfFileExists{xurl.sty}{\usepackage{xurl}}{} % add URL line breaks if available
\urlstyle{same}
\hypersetup{
  pdftitle={🚧Network Edges Inference🚧},
  colorlinks=true,
  linkcolor={blue},
  filecolor={Maroon},
  citecolor={Blue},
  urlcolor={Blue},
  pdfcreator={LaTeX via pandoc}}


\title{🚧Network Edges Inference🚧}
\author{}
\date{}
\begin{document}
\maketitle


\textless\textless\textless\textless\textless\textless\textless{} HEAD

This document provides a unified framework for performing Bayesian
inference of social networks based on observational data, implementing
the BISoN framework methodologies Hart et al. (2023). Real-world social
network connections (edges) are often inferred from observational data
where uniform sampling over all pairs is not always possible. This
framework enables modeling edge weights with uncertainty by decomposing
aggregated observations into the generative sampling processes.

Below, we detail the implementation of various edge weight models for
different types of observational data using the Bayesian Inference
(\textbf{BI}) package via Python, R, and Julia.

\section{\#\# General Principles}\label{general-principles}

\subsection{General Principles}\label{general-principles-1}

This document implements the BISoN framework methodologies Hart et al.
(2023) that enables modeling edge weights with uncertainty by
decomposing aggregated observations into the generative sampling
processes. Below, we detail the implementation of various edge weight
models for different types of observational data using the Bayesian
Inference (\textbf{BI}) package via Python, R, and Julia.

\begin{quote}
\begin{quote}
\begin{quote}
\begin{quote}
\begin{quote}
\begin{quote}
\begin{quote}
20c35a832b7f5f82f2011e8aad8fe0f1e2851c1b The fundamental principle
behind the Network Edges Inference models is to capture uncertainty in
edge weights by modeling the data collection process directly, rather
than relying on point estimates (e.g., Simple Ratio Index). By
decomposing the data into distinct observation periods, these models can
inherently capture the uncertainty generated by sampling effort and
allow for observation-level effects (such as variable visibility) to be
controlled.
\end{quote}
\end{quote}
\end{quote}
\end{quote}
\end{quote}
\end{quote}
\end{quote}

Three main types of observational data are supported: 1. \textbf{Binary
data}: Presence/absence of social events in fixed sampling periods. 2.
\textbf{Count data}: Number of social events per sampling period. 3.
\textbf{Duration data}: Amount of time spent engaged in social events
within a sampling period.

\subsection{\textless\textless\textless\textless\textless\textless\textless{}
HEAD}\label{head}

\subsubsection{Binary Data Model}\label{binary-data-model}

\section{\#\#\#\# General Principles}\label{general-principles-2}

For each models types Hart et al. (2023) provided the possibility to
account for: 1. \textbf{Directionality}: Directional or non-directional
edge weights can be modeled. 2. \textbf{Fixed effects}: Additional
variables accounting for fixed effects (e.g., location or age
difference) can be added to the predictor iteratively. 3. \textbf{Random
effects}: Additional variables accounting for nodal random effects
(e.g., location or age difference) can be added to the predictor
iteratively. 4. \textbf{Partial pooling} : Additional variables
accounting for dyads random effects (e.g., goup belonging) can be added
to the predictor iteratively. 5. \textbf{Zero-inflation}: Additional
variables accounting for random effects (e.g., location or age
difference) can be added to the predictor iteratively.

\begin{center}\rule{0.5\linewidth}{0.5pt}\end{center}

\subsection{Binary Data Model}\label{binary-data-model-1}

\subsubsection{General Principles}\label{general-principles-3}

\begin{quote}
\begin{quote}
\begin{quote}
\begin{quote}
\begin{quote}
\begin{quote}
\begin{quote}
20c35a832b7f5f82f2011e8aad8fe0f1e2851c1b The binary data model is used
when the presence or absence of a social event for each dyad is recorded
in each sampling period. The edge weight represents the probability of
(or proportion of time) engaging in a social event in a fixed period of
time.
\end{quote}
\end{quote}
\end{quote}
\end{quote}
\end{quote}
\end{quote}
\end{quote}

\paragraph{Considerations}\label{considerations}

\begin{tcolorbox}[enhanced jigsaw, colframe=quarto-callout-note-color-frame, bottomtitle=1mm, coltitle=black, titlerule=0mm, colbacktitle=quarto-callout-note-color!10!white, colback=white, opacitybacktitle=0.6, rightrule=.15mm, toptitle=1mm, opacityback=0, arc=.35mm, left=2mm, title=\textcolor{quarto-callout-note-color}{\faInfo}\hspace{0.5em}{Note}, breakable, bottomrule=.15mm, toprule=.15mm, leftrule=.75mm]

\begin{itemize}
\tightlist
\item
  \textbf{Priors}: The edge weight parameters normally represent
  logit-scale proportions and typically use standard Normal priors
  relative to prior structural beliefs (e.g., \texttt{Normal(0,\ 2)}).
\item
  \textbf{Data}: The model uses the number of sampling periods
  (\texttt{divisor}) where an event could occur and the number of
  observed events (\texttt{event}).
\item
  \textbf{Fixed Effects}: Additional variables accounting for fixed
  effects (e.g., location or age difference) can be added to the
  predictor iteratively.
\end{itemize}

\end{tcolorbox}

\textless\textless\textless\textless\textless\textless\textless{} HEAD
\#\#\#\# Example with simulated data ======= \#\#\# Example with
simulated data Simulated data can be generated with \textbf{bisonR}
function \texttt{simulate\_bison\_model}. For all example we will
account the presence of all possible features: 1 fixed nodal predictors,
1 random nodal predictors, presence of subgroups, partial\_pooling and
zero inflated model. This will allow you to adapt the model by removing
parts of the model based on your data.

To be compatible with \texttt{data} used in the model \texttt{data} must
contain and be zero-based indexed (i.e.~dyad\_idsID, num\_random\_groups
need to start at 0): 1. \emph{num\_rows} : Number of observations for
the total period of time. 2. \emph{event} : For binary model whether
event occur (1) or not (0). 3. \emph{duration} : Duration for each
observations. 4. \emph{dyad\_ids} : ID for each possible combinations of
dyad. 5. \emph{num\_edges}: For binary model this is equivalent of
\emph{num\_rows}. 6. \emph{num\_fixed}: Number of fixed nodal
predictors. 7. \emph{num\_random}: Number of random nodal predictors. 8.
\emph{num\_random\_groups} : Number of subgroups. 9.
\emph{random\_group\_index} : Subgroups ID. 10. \emph{design\_fixed}:
Design matrix (N, P) for nodal predictors, where N is the number of
individuals and P the number of predictors. 11. \emph{design\_random}:
Design matrix for nodal random effects. 12. \emph{partial\_pooling}:
Indicator (0 or 1) for enabling/disabling partial pooling of edge
weights. 13. \emph{zero\_inflated}: Indicator (0 or 1) for
enabling/disabling zero-inflation.

\begin{quote}
\begin{quote}
\begin{quote}
\begin{quote}
\begin{quote}
\begin{quote}
\begin{quote}
20c35a832b7f5f82f2011e8aad8fe0f1e2851c1b ::: \{.panel-tabset
group=``language''\} \#\# Python
\end{quote}
\end{quote}
\end{quote}
\end{quote}
\end{quote}
\end{quote}
\end{quote}

\begin{Shaded}
\begin{Highlighting}[]
\ImportTok{from}\NormalTok{ BI }\ImportTok{import}\NormalTok{ bi}

\NormalTok{m }\OperatorTok{=}\NormalTok{ bi(platform}\OperatorTok{=}\StringTok{\textquotesingle{}cpu\textquotesingle{}}\NormalTok{)}
\OperatorTok{\textless{}\textless{}\textless{}\textless{}\textless{}\textless{}\textless{}}\NormalTok{ HEAD}
\NormalTok{m.data\_on\_model }\OperatorTok{=}\NormalTok{ \{}\StringTok{\textquotesingle{}data\textquotesingle{}}\NormalTok{: data\_list\_jax\} }\CommentTok{\# Must contain \textquotesingle{}prior\_edge\_mu\textquotesingle{}, \textquotesingle{}num\_edges\textquotesingle{}, \textquotesingle{}dyad\_ids\textquotesingle{}, \textquotesingle{}divisor\textquotesingle{}, \textquotesingle{}event\textquotesingle{}, etc.}

\KeywordTok{def}\NormalTok{ bi\_model\_binary(data):}
    \CommentTok{\# Edge weights prior}
\NormalTok{    edge\_weight }\OperatorTok{=}\NormalTok{ m.dist.normal(data[}\StringTok{\textquotesingle{}prior\_edge\_mu\textquotesingle{}}\NormalTok{], data[}\StringTok{\textquotesingle{}prior\_edge\_sigma\textquotesingle{}}\NormalTok{], }
\NormalTok{                                shape}\OperatorTok{=}\NormalTok{(data[}\StringTok{\textquotesingle{}num\_edges\textquotesingle{}}\NormalTok{],), name}\OperatorTok{=}\StringTok{"edge\_weight"}\NormalTok{)}
    
    \CommentTok{\# 0{-}based indexing for JAX}
\NormalTok{    dyad\_ids\_0 }\OperatorTok{=}\NormalTok{ data[}\StringTok{\textquotesingle{}dyad\_ids\textquotesingle{}}\NormalTok{] }\OperatorTok{{-}} \DecValTok{1}
    
    \CommentTok{\# Build predictor}
\NormalTok{    predictor }\OperatorTok{=}\NormalTok{ edge\_weight[dyad\_ids\_0]}
    
    \CommentTok{\# Add fixed effects if applicable}
    \ControlFlowTok{if}\NormalTok{ data[}\StringTok{\textquotesingle{}num\_fixed\textquotesingle{}}\NormalTok{] }\OperatorTok{\textgreater{}} \DecValTok{0}\NormalTok{:}
\NormalTok{        beta\_fixed }\OperatorTok{=}\NormalTok{ m.dist.normal(data[}\StringTok{\textquotesingle{}prior\_fixed\_mu\textquotesingle{}}\NormalTok{], data[}\StringTok{\textquotesingle{}prior\_fixed\_sigma\textquotesingle{}}\NormalTok{], }
\NormalTok{                                   shape}\OperatorTok{=}\NormalTok{(data[}\StringTok{\textquotesingle{}num\_fixed\textquotesingle{}}\NormalTok{],), name}\OperatorTok{=}\StringTok{"beta\_fixed"}\NormalTok{)}
\NormalTok{        predictor }\OperatorTok{=}\NormalTok{ predictor }\OperatorTok{+}\NormalTok{ m.jnp.dot(data[}\StringTok{\textquotesingle{}design\_fixed\textquotesingle{}}\NormalTok{], beta\_fixed)}
    
    \CommentTok{\# Convert from logit scale to probability}
\NormalTok{    p }\OperatorTok{=}\NormalTok{ m.jax.scipy.special.expit(predictor)}
    
    \CommentTok{\# Binomial Likelihood}
\NormalTok{    m.dist.binomial(data[}\StringTok{\textquotesingle{}divisor\textquotesingle{}}\NormalTok{], p, obs}\OperatorTok{=}\NormalTok{data[}\StringTok{\textquotesingle{}event\textquotesingle{}}\NormalTok{], name}\OperatorTok{=}\StringTok{"event"}\NormalTok{)}

\NormalTok{m.fit(bi\_model\_binary)}
\NormalTok{m.summary()}
\end{Highlighting}
\end{Shaded}

\subsection{R}\label{r}

\begin{Shaded}
\begin{Highlighting}[]
\FunctionTok{library}\NormalTok{(BayesianInference)}
\NormalTok{m }\OtherTok{\textless{}{-}} \FunctionTok{importBI}\NormalTok{(}\AttributeTok{platform =} \StringTok{"cpu"}\NormalTok{)}
\NormalTok{jnp }\OtherTok{\textless{}{-}}\NormalTok{ reticulate}\SpecialCharTok{::}\FunctionTok{import}\NormalTok{(}\StringTok{"jax.numpy"}\NormalTok{)}
\NormalTok{jax\_scipy }\OtherTok{\textless{}{-}}\NormalTok{ reticulate}\SpecialCharTok{::}\FunctionTok{import}\NormalTok{(}\StringTok{"jax.scipy.special"}\NormalTok{)}

\NormalTok{m}\SpecialCharTok{$}\NormalTok{data\_on\_model }\OtherTok{\textless{}{-}} \FunctionTok{list}\NormalTok{(}\AttributeTok{data =}\NormalTok{ data\_list\_jax)}

\NormalTok{bi\_model\_binary }\OtherTok{\textless{}{-}} \ControlFlowTok{function}\NormalTok{(data) \{}
  \CommentTok{\# Edge weights prior}
\NormalTok{  edge\_weight }\OtherTok{\textless{}{-}}\NormalTok{ m}\SpecialCharTok{$}\NormalTok{dist}\SpecialCharTok{$}\FunctionTok{normal}\NormalTok{(data}\SpecialCharTok{$}\NormalTok{prior\_edge\_mu, data}\SpecialCharTok{$}\NormalTok{prior\_edge\_sigma, }
                               \AttributeTok{shape=}\FunctionTok{tuple}\NormalTok{(data}\SpecialCharTok{$}\NormalTok{num\_edges), }\AttributeTok{name=}\StringTok{"edge\_weight"}\NormalTok{)}
  
  \CommentTok{\# 0{-}based indexing for JAX}
\NormalTok{  dyad\_ids\_0 }\OtherTok{\textless{}{-}}\NormalTok{ data}\SpecialCharTok{$}\NormalTok{dyad\_ids }\SpecialCharTok{{-}} \DecValTok{1}\NormalTok{L}
  
  \CommentTok{\# Build predictor}
\NormalTok{  predictor }\OtherTok{\textless{}{-}}\NormalTok{ edge\_weight[dyad\_ids\_0]}
  
  \ControlFlowTok{if}\NormalTok{ (data}\SpecialCharTok{$}\NormalTok{num\_fixed }\SpecialCharTok{\textgreater{}} \DecValTok{0}\NormalTok{) \{}
\NormalTok{    beta\_fixed }\OtherTok{\textless{}{-}}\NormalTok{ m}\SpecialCharTok{$}\NormalTok{dist}\SpecialCharTok{$}\FunctionTok{normal}\NormalTok{(data}\SpecialCharTok{$}\NormalTok{prior\_fixed\_mu, data}\SpecialCharTok{$}\NormalTok{prior\_fixed\_sigma, }
                                \AttributeTok{shape=}\FunctionTok{tuple}\NormalTok{(data}\SpecialCharTok{$}\NormalTok{num\_fixed), }\AttributeTok{name=}\StringTok{"beta\_fixed"}\NormalTok{)}
\NormalTok{    predictor }\OtherTok{\textless{}{-}}\NormalTok{ predictor }\SpecialCharTok{+}\NormalTok{ jnp}\SpecialCharTok{$}\FunctionTok{dot}\NormalTok{(data}\SpecialCharTok{$}\NormalTok{design\_fixed, beta\_fixed)}
\NormalTok{  \}}
  
\NormalTok{  p }\OtherTok{\textless{}{-}}\NormalTok{ jax\_scipy}\SpecialCharTok{$}\FunctionTok{expit}\NormalTok{(predictor)}
  
\NormalTok{  m}\SpecialCharTok{$}\NormalTok{dist}\SpecialCharTok{$}\FunctionTok{binomial}\NormalTok{(data}\SpecialCharTok{$}\NormalTok{divisor, p, }\AttributeTok{obs=}\NormalTok{data}\SpecialCharTok{$}\NormalTok{event, }\AttributeTok{name=}\StringTok{"event"}\NormalTok{)}
\NormalTok{\}}

\NormalTok{m}\SpecialCharTok{$}\FunctionTok{fit}\NormalTok{(bi\_model\_binary)}
\NormalTok{m}\SpecialCharTok{$}\FunctionTok{summary}\NormalTok{()}
\end{Highlighting}
\end{Shaded}

\subsection{Julia}\label{julia}

\begin{Shaded}
\begin{Highlighting}[]
\ImportTok{using} \BuiltInTok{BayesianInference}

\NormalTok{m }\OperatorTok{=} \FunctionTok{importBI}\NormalTok{(platform}\OperatorTok{=}\StringTok{"cpu"}\NormalTok{)}
\NormalTok{m.data\_on\_model }\OperatorTok{=} \FunctionTok{Dict}\NormalTok{(}\StringTok{"data"} \OperatorTok{=\textgreater{}}\NormalTok{ data\_list\_jax)}

\PreprocessorTok{@BI} \KeywordTok{function} \FunctionTok{bi\_model\_binary}\NormalTok{(data)}
\NormalTok{    edge\_weight }\OperatorTok{=}\NormalTok{ m.dist.}\FunctionTok{normal}\NormalTok{(data[}\StringTok{"prior\_edge\_mu"}\NormalTok{], data[}\StringTok{"prior\_edge\_sigma"}\NormalTok{], }
\NormalTok{                                shape}\OperatorTok{=}\NormalTok{(data[}\StringTok{"num\_edges"}\NormalTok{],), name}\OperatorTok{=}\StringTok{"edge\_weight"}\NormalTok{)}
    
    \CommentTok{\# 0{-}based indexing for JAX inside the environment}
\NormalTok{    dyad\_ids\_0 }\OperatorTok{=}\NormalTok{ data[}\StringTok{"dyad\_ids"}\NormalTok{] }\OperatorTok{.{-}} \FloatTok{1}
    
\NormalTok{    predictor }\OperatorTok{=}\NormalTok{ edge\_weight[dyad\_ids\_0]}
    
    \ControlFlowTok{if}\NormalTok{ data[}\StringTok{"num\_fixed"}\NormalTok{] }\OperatorTok{\textgreater{}} \FloatTok{0}
\NormalTok{        beta\_fixed }\OperatorTok{=}\NormalTok{ m.dist.}\FunctionTok{normal}\NormalTok{(data[}\StringTok{"prior\_fixed\_mu"}\NormalTok{], data[}\StringTok{"prior\_fixed\_sigma"}\NormalTok{], }
\NormalTok{                                   shape}\OperatorTok{=}\NormalTok{(data[}\StringTok{"num\_fixed"}\NormalTok{],), name}\OperatorTok{=}\StringTok{"beta\_fixed"}\NormalTok{)}
\NormalTok{        predictor }\OperatorTok{=}\NormalTok{ predictor }\OperatorTok{.+}\NormalTok{ jnp.}\FunctionTok{dot}\NormalTok{(data[}\StringTok{"design\_fixed"}\NormalTok{], beta\_fixed)}
    \ControlFlowTok{end}
    
\NormalTok{    p }\OperatorTok{=}\NormalTok{ jax.scipy.special.}\FunctionTok{expit}\NormalTok{(predictor)}
    
\NormalTok{    m.dist.}\FunctionTok{binomial}\NormalTok{(data[}\StringTok{"divisor"}\NormalTok{], p, obs}\OperatorTok{=}\NormalTok{data[}\StringTok{"event"}\NormalTok{], name}\OperatorTok{=}\StringTok{"event"}\NormalTok{)}
\KeywordTok{end}

\NormalTok{m.}\FunctionTok{fit}\NormalTok{(bi\_model\_binary)}
\NormalTok{m.}\FunctionTok{summary}\NormalTok{()}
\end{Highlighting}
\end{Shaded}

:::

\paragraph{Mathematical Details}\label{mathematical-details}

For binary data, the mathematical process applies a Binomial likelihood
assuming multiple Bernoulli trials: \[
X_{ij}^{(n)} \sim \text{Binomial}(D_{ij}^{(n)}, p_{ij}^{(n)})
\] Where \(X_{ij}^{(n)}\) is the presence/absence in the \(n\)-th
observation, \(D_{ij}^{(n)}\) is the condition (often 1 or representing
multiple grouped trials), and \(p_{ij}^{(n)}\) is the dyad interaction
probability. The model logit-links the predictor effects with the
primary edge weight: \[
\text{logit}(p_{ij}^{(n)}) = \omega_{ij} + \dots
\] Where \(\omega_{ij}\) represents the foundational edge weight
parameter for dyad \(i, j\).

\paragraph{Notes}\label{notes}

\begin{tcolorbox}[enhanced jigsaw, colframe=quarto-callout-note-color-frame, bottomtitle=1mm, coltitle=black, titlerule=0mm, colbacktitle=quarto-callout-note-color!10!white, colback=white, opacitybacktitle=0.6, rightrule=.15mm, toptitle=1mm, opacityback=0, arc=.35mm, left=2mm, title=\textcolor{quarto-callout-note-color}{\faInfo}\hspace{0.5em}{Note}, breakable, bottomrule=.15mm, toprule=.15mm, leftrule=.75mm]

The binary binomial implementation directly provides a framework
mirroring the Simple Ratio Index (SRI) while preserving variance from
low sample depth.

\end{tcolorbox}

\begin{center}\rule{0.5\linewidth}{0.5pt}\end{center}

\subsubsection{Count Data Model}\label{count-data-model}

\paragraph{General Principles}\label{general-principles-4}

The count data model interprets instances where the exact \emph{number
of events} for each dyad is recorded in each sampling period. Under this
formulation, the model outputs a rate of occurrence of social events per
unit of time instead of a bounded probability.

\paragraph{Considerations}\label{considerations-1}

\begin{tcolorbox}[enhanced jigsaw, colframe=quarto-callout-note-color-frame, bottomtitle=1mm, coltitle=black, titlerule=0mm, colbacktitle=quarto-callout-note-color!10!white, colback=white, opacitybacktitle=0.6, rightrule=.15mm, toptitle=1mm, opacityback=0, arc=.35mm, left=2mm, title=\textcolor{quarto-callout-note-color}{\faInfo}\hspace{0.5em}{Note}, breakable, bottomrule=.15mm, toprule=.15mm, leftrule=.75mm]

\begin{itemize}
\tightlist
\item
  \textbf{Priors}: Counts use a standard Poisson process; therefore,
  edge predictors dictate the log-rate (\(\lambda\)) and must be
  carefully scaled relative to time. \texttt{Normal} priors on the
  log-rate parameter effectively capture the magnitudes.
\item
  \textbf{Link Function}: Uses the exponential function applied to the
  linear predictor configuration.
\end{itemize}

\end{tcolorbox}

\paragraph{Example with simulated
data}\label{example-with-simulated-data}

\subsection{Python}

\begin{Shaded}
\begin{Highlighting}[]
\ImportTok{from}\NormalTok{ BI }\ImportTok{import}\NormalTok{ bi}

\NormalTok{m }\OperatorTok{=}\NormalTok{ bi(platform}\OperatorTok{=}\StringTok{\textquotesingle{}cpu\textquotesingle{}}\NormalTok{)}
\NormalTok{m.data\_on\_model }\OperatorTok{=}\NormalTok{ \{}\StringTok{\textquotesingle{}data\textquotesingle{}}\NormalTok{: data\_list\_jax\} }

\KeywordTok{def}\NormalTok{ bi\_model\_count(data):}
    \CommentTok{\# Edge weights prior (log{-}rate)}
\NormalTok{    edge\_weight }\OperatorTok{=}\NormalTok{ m.dist.normal(data[}\StringTok{\textquotesingle{}prior\_edge\_mu\textquotesingle{}}\NormalTok{], data[}\StringTok{\textquotesingle{}prior\_edge\_sigma\textquotesingle{}}\NormalTok{], }
\NormalTok{                                shape}\OperatorTok{=}\NormalTok{(data[}\StringTok{\textquotesingle{}num\_edges\textquotesingle{}}\NormalTok{],), name}\OperatorTok{=}\StringTok{"edge\_weight"}\NormalTok{)}
    
\NormalTok{    dyad\_ids\_0 }\OperatorTok{=}\NormalTok{ data[}\StringTok{\textquotesingle{}dyad\_ids\textquotesingle{}}\NormalTok{] }\OperatorTok{{-}} \DecValTok{1}
\NormalTok{    predictor }\OperatorTok{=}\NormalTok{ edge\_weight[dyad\_ids\_0]}
    
    \ControlFlowTok{if}\NormalTok{ data[}\StringTok{\textquotesingle{}num\_fixed\textquotesingle{}}\NormalTok{] }\OperatorTok{\textgreater{}} \DecValTok{0}\NormalTok{:}
\NormalTok{        beta\_fixed }\OperatorTok{=}\NormalTok{ m.dist.normal(data[}\StringTok{\textquotesingle{}prior\_fixed\_mu\textquotesingle{}}\NormalTok{], data[}\StringTok{\textquotesingle{}prior\_fixed\_sigma\textquotesingle{}}\NormalTok{], }
\NormalTok{                                   shape}\OperatorTok{=}\NormalTok{(data[}\StringTok{\textquotesingle{}num\_fixed\textquotesingle{}}\NormalTok{],), name}\OperatorTok{=}\StringTok{"beta\_fixed"}\NormalTok{)}
\NormalTok{        predictor }\OperatorTok{=}\NormalTok{ predictor }\OperatorTok{+}\NormalTok{ m.jnp.dot(data[}\StringTok{\textquotesingle{}design\_fixed\textquotesingle{}}\NormalTok{], beta\_fixed)}
    
    \CommentTok{\# Scale via event lengths/time blocks}
\NormalTok{    rate }\OperatorTok{=}\NormalTok{ m.jnp.exp(predictor) }\OperatorTok{*}\NormalTok{ data[}\StringTok{\textquotesingle{}divisor\textquotesingle{}}\NormalTok{]}
    
    \CommentTok{\# Poisson Likelihood}
\NormalTok{    m.dist.poisson(rate, obs}\OperatorTok{=}\NormalTok{data[}\StringTok{\textquotesingle{}event\textquotesingle{}}\NormalTok{], name}\OperatorTok{=}\StringTok{"event"}\NormalTok{)}

\NormalTok{m.fit(bi\_model\_count)}
\NormalTok{m.summary()}
\end{Highlighting}
\end{Shaded}

\subsection{R}

\begin{Shaded}
\begin{Highlighting}[]
\FunctionTok{library}\NormalTok{(BayesianInference)}
\NormalTok{m }\OtherTok{\textless{}{-}} \FunctionTok{importBI}\NormalTok{(}\AttributeTok{platform =} \StringTok{"cpu"}\NormalTok{)}
\NormalTok{jnp }\OtherTok{\textless{}{-}}\NormalTok{ reticulate}\SpecialCharTok{::}\FunctionTok{import}\NormalTok{(}\StringTok{"jax.numpy"}\NormalTok{)}

\NormalTok{m}\SpecialCharTok{$}\NormalTok{data\_on\_model }\OtherTok{\textless{}{-}} \FunctionTok{list}\NormalTok{(}\AttributeTok{data =}\NormalTok{ data\_list\_jax)}

\NormalTok{bi\_model\_count }\OtherTok{\textless{}{-}} \ControlFlowTok{function}\NormalTok{(data) \{}
\NormalTok{  edge\_weight }\OtherTok{\textless{}{-}}\NormalTok{ m}\SpecialCharTok{$}\NormalTok{dist}\SpecialCharTok{$}\FunctionTok{normal}\NormalTok{(data}\SpecialCharTok{$}\NormalTok{prior\_edge\_mu, data}\SpecialCharTok{$}\NormalTok{prior\_edge\_sigma, }
                               \AttributeTok{shape=}\FunctionTok{tuple}\NormalTok{(data}\SpecialCharTok{$}\NormalTok{num\_edges), }\AttributeTok{name=}\StringTok{"edge\_weight"}\NormalTok{)}
  
\NormalTok{  dyad\_ids\_0 }\OtherTok{\textless{}{-}}\NormalTok{ data}\SpecialCharTok{$}\NormalTok{dyad\_ids }\SpecialCharTok{{-}} \DecValTok{1}\NormalTok{L}
\NormalTok{  predictor }\OtherTok{\textless{}{-}}\NormalTok{ edge\_weight[dyad\_ids\_0]}
  
  \ControlFlowTok{if}\NormalTok{ (data}\SpecialCharTok{$}\NormalTok{num\_fixed }\SpecialCharTok{\textgreater{}} \DecValTok{0}\NormalTok{) \{}
\NormalTok{    beta\_fixed }\OtherTok{\textless{}{-}}\NormalTok{ m}\SpecialCharTok{$}\NormalTok{dist}\SpecialCharTok{$}\FunctionTok{normal}\NormalTok{(data}\SpecialCharTok{$}\NormalTok{prior\_fixed\_mu, data}\SpecialCharTok{$}\NormalTok{prior\_fixed\_sigma, }
                                \AttributeTok{shape=}\FunctionTok{tuple}\NormalTok{(data}\SpecialCharTok{$}\NormalTok{num\_fixed), }\AttributeTok{name=}\StringTok{"beta\_fixed"}\NormalTok{)}
\NormalTok{    predictor }\OtherTok{\textless{}{-}}\NormalTok{ predictor }\SpecialCharTok{+}\NormalTok{ jnp}\SpecialCharTok{$}\FunctionTok{dot}\NormalTok{(data}\SpecialCharTok{$}\NormalTok{design\_fixed, beta\_fixed)}
\NormalTok{  \}}
  
\NormalTok{  rate }\OtherTok{\textless{}{-}}\NormalTok{ jnp}\SpecialCharTok{$}\FunctionTok{exp}\NormalTok{(predictor) }\SpecialCharTok{*}\NormalTok{ data}\SpecialCharTok{$}\NormalTok{divisor}
  
\NormalTok{  m}\SpecialCharTok{$}\NormalTok{dist}\SpecialCharTok{$}\FunctionTok{poisson}\NormalTok{(rate, }\AttributeTok{obs=}\NormalTok{data}\SpecialCharTok{$}\NormalTok{event, }\AttributeTok{name=}\StringTok{"event"}\NormalTok{)}
\NormalTok{\}}

\NormalTok{m}\SpecialCharTok{$}\FunctionTok{fit}\NormalTok{(bi\_model\_count)}
\NormalTok{m}\SpecialCharTok{$}\FunctionTok{summary}\NormalTok{()}
\end{Highlighting}
\end{Shaded}

\subsection{Julia}

\begin{Shaded}
\begin{Highlighting}[]
\ImportTok{using} \BuiltInTok{BayesianInference}

\NormalTok{m }\OperatorTok{=} \FunctionTok{importBI}\NormalTok{(platform}\OperatorTok{=}\StringTok{"cpu"}\NormalTok{)}
\NormalTok{m.data\_on\_model }\OperatorTok{=} \FunctionTok{Dict}\NormalTok{(}\StringTok{"data"} \OperatorTok{=\textgreater{}}\NormalTok{ data\_list\_jax)}

\PreprocessorTok{@BI} \KeywordTok{function} \FunctionTok{bi\_model\_count}\NormalTok{(data)}
\NormalTok{    edge\_weight }\OperatorTok{=}\NormalTok{ m.dist.}\FunctionTok{normal}\NormalTok{(data[}\StringTok{"prior\_edge\_mu"}\NormalTok{], data[}\StringTok{"prior\_edge\_sigma"}\NormalTok{], }
\NormalTok{                                shape}\OperatorTok{=}\NormalTok{(data[}\StringTok{"num\_edges"}\NormalTok{],), name}\OperatorTok{=}\StringTok{"edge\_weight"}\NormalTok{)}
    
\NormalTok{    dyad\_ids\_0 }\OperatorTok{=}\NormalTok{ data[}\StringTok{"dyad\_ids"}\NormalTok{] }\OperatorTok{.{-}} \FloatTok{1}
\NormalTok{    predictor }\OperatorTok{=}\NormalTok{ edge\_weight[dyad\_ids\_0]}
    
    \ControlFlowTok{if}\NormalTok{ data[}\StringTok{"num\_fixed"}\NormalTok{] }\OperatorTok{\textgreater{}} \FloatTok{0}
\NormalTok{        beta\_fixed }\OperatorTok{=}\NormalTok{ m.dist.}\FunctionTok{normal}\NormalTok{(data[}\StringTok{"prior\_fixed\_mu"}\NormalTok{], data[}\StringTok{"prior\_fixed\_sigma"}\NormalTok{], }
\NormalTok{                                   shape}\OperatorTok{=}\NormalTok{(data[}\StringTok{"num\_fixed"}\NormalTok{],), name}\OperatorTok{=}\StringTok{"beta\_fixed"}\NormalTok{)}
\NormalTok{        predictor }\OperatorTok{=}\NormalTok{ predictor }\OperatorTok{.+}\NormalTok{ jnp.}\FunctionTok{dot}\NormalTok{(data[}\StringTok{"design\_fixed"}\NormalTok{], beta\_fixed)}
    \ControlFlowTok{end}
    
\NormalTok{    rate }\OperatorTok{=}\NormalTok{ jnp.}\FunctionTok{exp}\NormalTok{(predictor) }\OperatorTok{.*}\NormalTok{ data[}\StringTok{"divisor"}\NormalTok{]}
    
\NormalTok{    m.dist.}\FunctionTok{poisson}\NormalTok{(rate, obs}\OperatorTok{=}\NormalTok{data[}\StringTok{"event"}\NormalTok{], name}\OperatorTok{=}\StringTok{"event"}\NormalTok{)}
\KeywordTok{end}

\NormalTok{m.}\FunctionTok{fit}\NormalTok{(bi\_model\_count)}
\NormalTok{m.}\FunctionTok{summary}\NormalTok{()}
\end{Highlighting}
\end{Shaded}

\paragraph{Mathematical Details}\label{mathematical-details-1}

For count data, we use the sum of Poisson-distributed random variables
to define interactions scaling with observation duration. \[
X_{ij}^{(n)} \sim \text{Poisson}(\lambda_{ij}^{(n)} D_{ij}^{(n)})
\] Where the overall rate modifier is described logarithmically: \[
\log(\lambda_{ij}^{(n)}) = \omega_{ij} + \dots
\] Where \(\omega_{ij}\) evaluates to the maximum likelihood estimation
counterpart of straightforward event rate counts across cumulative time
blocks.

\paragraph{Notes}\label{notes-1}

\begin{tcolorbox}[enhanced jigsaw, colframe=quarto-callout-note-color-frame, bottomtitle=1mm, coltitle=black, titlerule=0mm, colbacktitle=quarto-callout-note-color!10!white, colback=white, opacitybacktitle=0.6, rightrule=.15mm, toptitle=1mm, opacityback=0, arc=.35mm, left=2mm, title=\textcolor{quarto-callout-note-color}{\faInfo}\hspace{0.5em}{Note}, breakable, bottomrule=.15mm, toprule=.15mm, leftrule=.75mm]

The Poisson distribution inherently manages the aggregation mechanics
allowing the edge weight to serve dynamically scaled estimates
independent from disparate observation lengths.

\end{tcolorbox}

\begin{center}\rule{0.5\linewidth}{0.5pt}\end{center}

\subsubsection{Duration Data Model}\label{duration-data-model}

\paragraph{General Principles}\label{general-principles-5}

The duration model simultaneously handles two data metrics: the
\emph{duration} of social events and the \emph{frequency} with which
they occur. It treats total engaged time mathematically using a
composition approach tracking exponential duration periods bounded
across a baseline Poisson event initiation rate.

\paragraph{Considerations}\label{considerations-2}

\begin{tcolorbox}[enhanced jigsaw, colframe=quarto-callout-note-color-frame, bottomtitle=1mm, coltitle=black, titlerule=0mm, colbacktitle=quarto-callout-note-color!10!white, colback=white, opacitybacktitle=0.6, rightrule=.15mm, toptitle=1mm, opacityback=0, arc=.35mm, left=2mm, title=\textcolor{quarto-callout-note-color}{\faInfo}\hspace{0.5em}{Note}, breakable, bottomrule=.15mm, toprule=.15mm, leftrule=.75mm]

\begin{itemize}
\tightlist
\item
  \textbf{Priors}: Requires \emph{two} primary priors. The standard
  edge-weight proportion models the relative event time (\(\omega\)),
  while an overarching generalized \texttt{rate} component drives the
  base observation scale.
\item
  \textbf{Complexity}: Fits two separate likelihood distributions
  (Poisson and Exponential) interacting across shared predictors to map
  both count frequencies and continuous lengths simultaneously.
\end{itemize}

\end{tcolorbox}

\paragraph{Example with simulated
data}\label{example-with-simulated-data-1}

\subsection{Python}

\begin{Shaded}
\begin{Highlighting}[]
\ImportTok{from}\NormalTok{ BI }\ImportTok{import}\NormalTok{ bi}

\NormalTok{m }\OperatorTok{=}\NormalTok{ bi(platform}\OperatorTok{=}\StringTok{\textquotesingle{}cpu\textquotesingle{}}\NormalTok{)}
\NormalTok{m.data\_on\_model }\OperatorTok{=}\NormalTok{ \{}\StringTok{\textquotesingle{}data\textquotesingle{}}\NormalTok{: data\_list\_jax\} }

\KeywordTok{def}\NormalTok{ bi\_model\_duration(data):}
    \CommentTok{\# Edge weights (proportion scale mapping)}
\NormalTok{    edge\_weight }\OperatorTok{=}\NormalTok{ m.dist.normal(data[}\StringTok{\textquotesingle{}prior\_edge\_mu\textquotesingle{}}\NormalTok{], data[}\StringTok{\textquotesingle{}prior\_edge\_sigma\textquotesingle{}}\NormalTok{], }
\NormalTok{                                shape}\OperatorTok{=}\NormalTok{(data[}\StringTok{\textquotesingle{}num\_edges\textquotesingle{}}\NormalTok{],), name}\OperatorTok{=}\StringTok{"edge\_weight"}\NormalTok{)}
    
    \CommentTok{\# Additional absolute rate parameter strictly positive (Half{-}Normal equivalent constraint logic)}
\NormalTok{    rate\_param }\OperatorTok{=}\NormalTok{ m.dist.normal(}\FloatTok{0.0}\NormalTok{, data[}\StringTok{\textquotesingle{}prior\_rate\_sigma\textquotesingle{}}\NormalTok{], }
\NormalTok{                               shape}\OperatorTok{=}\NormalTok{(data[}\StringTok{\textquotesingle{}num\_edges\textquotesingle{}}\NormalTok{],), name}\OperatorTok{=}\StringTok{"rate"}\NormalTok{)}
\NormalTok{    rate\_positive }\OperatorTok{=}\NormalTok{ m.jnp.}\BuiltInTok{abs}\NormalTok{(rate\_param) }
    
\NormalTok{    dyad\_ids\_0 }\OperatorTok{=}\NormalTok{ data[}\StringTok{\textquotesingle{}dyad\_ids\textquotesingle{}}\NormalTok{] }\OperatorTok{{-}} \DecValTok{1}
\NormalTok{    predictor }\OperatorTok{=}\NormalTok{ edge\_weight[dyad\_ids\_0]}
    
    \ControlFlowTok{if}\NormalTok{ data[}\StringTok{\textquotesingle{}num\_fixed\textquotesingle{}}\NormalTok{] }\OperatorTok{\textgreater{}} \DecValTok{0}\NormalTok{:}
\NormalTok{        beta\_fixed }\OperatorTok{=}\NormalTok{ m.dist.normal(data[}\StringTok{\textquotesingle{}prior\_fixed\_mu\textquotesingle{}}\NormalTok{], data[}\StringTok{\textquotesingle{}prior\_fixed\_sigma\textquotesingle{}}\NormalTok{], }
\NormalTok{                                   shape}\OperatorTok{=}\NormalTok{(data[}\StringTok{\textquotesingle{}num\_fixed\textquotesingle{}}\NormalTok{],), name}\OperatorTok{=}\StringTok{"beta\_fixed"}\NormalTok{)}
\NormalTok{        predictor }\OperatorTok{=}\NormalTok{ predictor }\OperatorTok{+}\NormalTok{ m.jnp.dot(data[}\StringTok{\textquotesingle{}design\_fixed\textquotesingle{}}\NormalTok{], beta\_fixed)}

    \CommentTok{\# Convert to expected total scale fraction}
\NormalTok{    p }\OperatorTok{=}\NormalTok{ m.jax.scipy.special.expit(predictor)}
    
    \CommentTok{\# Build linked parameters}
\NormalTok{    lambda\_exp }\OperatorTok{=}\NormalTok{ rate\_positive[dyad\_ids\_0] }\OperatorTok{/}\NormalTok{ p}
\NormalTok{    lambda\_pois }\OperatorTok{=}\NormalTok{ rate\_positive }\OperatorTok{*}\NormalTok{ data[}\StringTok{\textquotesingle{}divisor\textquotesingle{}}\NormalTok{]}
    
    \CommentTok{\# Joint likelihood inference}
\OperatorTok{=======}
\NormalTok{m.data\_on\_model }\OperatorTok{=}\NormalTok{ \{}\StringTok{\textquotesingle{}data\textquotesingle{}}\NormalTok{: data\} }

\KeywordTok{def}\NormalTok{ bi\_model\_duration(data):}

    \CommentTok{\# 1. Edge weights prior with partial pooling : Set the baseline (Intercept)}
\NormalTok{    edge\_sigma }\OperatorTok{=}\NormalTok{ m.dist.half\_normal(}\FloatTok{2.0}\NormalTok{, name}\OperatorTok{=}\StringTok{"edge\_sigma"}\NormalTok{)}
\NormalTok{    edge\_weight }\OperatorTok{=}\NormalTok{ m.dist.normal(}\FloatTok{0.0}\NormalTok{, edge\_sigma, shape}\OperatorTok{=}\NormalTok{(data[}\StringTok{\textquotesingle{}num\_edges\textquotesingle{}}\NormalTok{],), name}\OperatorTok{=}\StringTok{"edge\_weight"}\NormalTok{)}
    
    \CommentTok{\# Additional absolute rate parameter strictly positive}
\NormalTok{    rate\_positive }\OperatorTok{=}\NormalTok{ m.dist.half\_normal(}\FloatTok{1.0}\NormalTok{, shape}\OperatorTok{=}\NormalTok{(data[}\StringTok{\textquotesingle{}num\_edges\textquotesingle{}}\NormalTok{],), name}\OperatorTok{=}\StringTok{"rate"}\NormalTok{)}
    
\NormalTok{    dyad\_ids\_0 }\OperatorTok{=}\NormalTok{ data[}\StringTok{\textquotesingle{}dyad\_ids\textquotesingle{}}\NormalTok{] }\OperatorTok{{-}} \DecValTok{1}
    \CommentTok{\#\# Build predictor}
\NormalTok{    predictor }\OperatorTok{=}\NormalTok{ edge\_weight[dyad\_ids\_0]}
    
    \CommentTok{\# 2. Fixed nodal factor(s)}
\NormalTok{    beta\_fixed }\OperatorTok{=}\NormalTok{ m.dist.normal(}\FloatTok{0.0}\NormalTok{, }\FloatTok{1.0}\NormalTok{, shape}\OperatorTok{=}\NormalTok{(data[}\StringTok{\textquotesingle{}num\_fixed\textquotesingle{}}\NormalTok{],), name}\OperatorTok{=}\StringTok{"beta\_fixed"}\NormalTok{)}
\NormalTok{    predictor }\OperatorTok{=}\NormalTok{ predictor }\OperatorTok{+}\NormalTok{ m.jnp.dot(data[}\StringTok{\textquotesingle{}design\_fixed\textquotesingle{}}\NormalTok{], beta\_fixed)}

    \CommentTok{\# 3. Random nodal effect(s)}
\NormalTok{    random\_group\_mu }\OperatorTok{=}\NormalTok{ m.dist.normal(}\FloatTok{0.0}\NormalTok{, }\FloatTok{1.0}\NormalTok{, shape}\OperatorTok{=}\NormalTok{(data[}\StringTok{\textquotesingle{}num\_random\_groups\textquotesingle{}}\NormalTok{],), name}\OperatorTok{=}\StringTok{"random\_group\_mu"}\NormalTok{)}
\NormalTok{    random\_group\_sigma }\OperatorTok{=}\NormalTok{ m.dist.half\_normal(}\FloatTok{1.0}\NormalTok{, shape}\OperatorTok{=}\NormalTok{(data[}\StringTok{\textquotesingle{}num\_random\_groups\textquotesingle{}}\NormalTok{],), name}\OperatorTok{=}\StringTok{"random\_group\_sigma"}\NormalTok{)}
    
\NormalTok{    group\_idx\_0 }\OperatorTok{=}\NormalTok{ data[}\StringTok{\textquotesingle{}random\_group\_index\textquotesingle{}}\NormalTok{] }\OperatorTok{{-}} \DecValTok{1}
\NormalTok{    beta\_random }\OperatorTok{=}\NormalTok{ m.dist.normal(random\_group\_mu[group\_idx\_0], random\_group\_sigma[group\_idx\_0], shape}\OperatorTok{=}\NormalTok{(data[}\StringTok{\textquotesingle{}num\_random\textquotesingle{}}\NormalTok{],), name}\OperatorTok{=}\StringTok{"beta\_random"}\NormalTok{)}
    
\NormalTok{    predictor }\OperatorTok{=}\NormalTok{ predictor }\OperatorTok{+}\NormalTok{ m.jnp.dot(data[}\StringTok{\textquotesingle{}design\_random\textquotesingle{}}\NormalTok{], beta\_random)}

    \CommentTok{\# 4. Link function}
\NormalTok{    p }\OperatorTok{=}\NormalTok{ m.link.inv\_logit(predictor)}
    
\NormalTok{    lambda\_exp }\OperatorTok{=}\NormalTok{ rate\_positive[dyad\_ids\_0] }\OperatorTok{/}\NormalTok{ p}
\NormalTok{    lambda\_pois }\OperatorTok{=}\NormalTok{ rate\_positive }\OperatorTok{*}\NormalTok{ data[}\StringTok{\textquotesingle{}divisor\textquotesingle{}}\NormalTok{]}
    
    \CommentTok{\# 5. Joint likelihood inference}
\OperatorTok{\textgreater{}\textgreater{}\textgreater{}\textgreater{}\textgreater{}\textgreater{}\textgreater{}} \DecValTok{20}\ErrorTok{c35a832b7f5f82f2011e8aad8fe0f1e2851c1b}
\NormalTok{    m.dist.exponential(lambda\_exp, obs}\OperatorTok{=}\NormalTok{data[}\StringTok{\textquotesingle{}event\textquotesingle{}}\NormalTok{], name}\OperatorTok{=}\StringTok{"event"}\NormalTok{)}
\NormalTok{    m.dist.poisson(lambda\_pois, obs}\OperatorTok{=}\NormalTok{data[}\StringTok{\textquotesingle{}event\_count\textquotesingle{}}\NormalTok{], name}\OperatorTok{=}\StringTok{"event\_count"}\NormalTok{)}

\NormalTok{m.fit(bi\_model\_duration)}
\NormalTok{m.summary()}
\end{Highlighting}
\end{Shaded}

\subsection{R}

\begin{Shaded}
\begin{Highlighting}[]
\SpecialCharTok{\textless{}}\ErrorTok{\textless{}\textless{}\textless{}\textless{}\textless{}\textless{}}\NormalTok{ HEAD}
\FunctionTok{library}\NormalTok{(BayesianInference)}
\NormalTok{m }\OtherTok{\textless{}{-}} \FunctionTok{importBI}\NormalTok{(}\AttributeTok{platform =} \StringTok{"cpu"}\NormalTok{)}
\NormalTok{jnp }\OtherTok{\textless{}{-}}\NormalTok{ reticulate}\SpecialCharTok{::}\FunctionTok{import}\NormalTok{(}\StringTok{"jax.numpy"}\NormalTok{)}
\NormalTok{jax\_scipy }\OtherTok{\textless{}{-}}\NormalTok{ reticulate}\SpecialCharTok{::}\FunctionTok{import}\NormalTok{(}\StringTok{"jax.scipy.special"}\NormalTok{)}

\NormalTok{m}\SpecialCharTok{$}\NormalTok{data\_on\_model }\OtherTok{\textless{}{-}} \FunctionTok{list}\NormalTok{(}\AttributeTok{data =}\NormalTok{ data\_list\_jax)}

\NormalTok{bi\_model\_duration }\OtherTok{\textless{}{-}} \ControlFlowTok{function}\NormalTok{(data) \{}
\NormalTok{  edge\_weight }\OtherTok{\textless{}{-}}\NormalTok{ m}\SpecialCharTok{$}\NormalTok{dist}\SpecialCharTok{$}\FunctionTok{normal}\NormalTok{(data}\SpecialCharTok{$}\NormalTok{prior\_edge\_mu, data}\SpecialCharTok{$}\NormalTok{prior\_edge\_sigma, }
                               \AttributeTok{shape=}\FunctionTok{tuple}\NormalTok{(data}\SpecialCharTok{$}\NormalTok{num\_edges), }\AttributeTok{name=}\StringTok{"edge\_weight"}\NormalTok{)}
  
\NormalTok{  rate\_param }\OtherTok{\textless{}{-}}\NormalTok{ m}\SpecialCharTok{$}\NormalTok{dist}\SpecialCharTok{$}\FunctionTok{normal}\NormalTok{(}\FloatTok{0.0}\NormalTok{, data}\SpecialCharTok{$}\NormalTok{prior\_rate\_sigma, }
                              \AttributeTok{shape=}\FunctionTok{tuple}\NormalTok{(data}\SpecialCharTok{$}\NormalTok{num\_edges), }\AttributeTok{name=}\StringTok{"rate"}\NormalTok{)}
\NormalTok{  rate\_positive }\OtherTok{\textless{}{-}}\NormalTok{ jnp}\SpecialCharTok{$}\FunctionTok{abs}\NormalTok{(rate\_param)}
  
\NormalTok{  dyad\_ids\_0 }\OtherTok{\textless{}{-}}\NormalTok{ data}\SpecialCharTok{$}\NormalTok{dyad\_ids }\SpecialCharTok{{-}} \DecValTok{1}\NormalTok{L}
\NormalTok{  predictor }\OtherTok{\textless{}{-}}\NormalTok{ edge\_weight[dyad\_ids\_0]}
  
  \ControlFlowTok{if}\NormalTok{ (data}\SpecialCharTok{$}\NormalTok{num\_fixed }\SpecialCharTok{\textgreater{}} \DecValTok{0}\NormalTok{) \{}
\NormalTok{    beta\_fixed }\OtherTok{\textless{}{-}}\NormalTok{ m}\SpecialCharTok{$}\NormalTok{dist}\SpecialCharTok{$}\FunctionTok{normal}\NormalTok{(data}\SpecialCharTok{$}\NormalTok{prior\_fixed\_mu, data}\SpecialCharTok{$}\NormalTok{prior\_fixed\_sigma, }
                                \AttributeTok{shape=}\FunctionTok{tuple}\NormalTok{(data}\SpecialCharTok{$}\NormalTok{num\_fixed), }\AttributeTok{name=}\StringTok{"beta\_fixed"}\NormalTok{)}
\NormalTok{    predictor }\OtherTok{\textless{}{-}}\NormalTok{ predictor }\SpecialCharTok{+}\NormalTok{ jnp}\SpecialCharTok{$}\FunctionTok{dot}\NormalTok{(data}\SpecialCharTok{$}\NormalTok{design\_fixed, beta\_fixed)}
\NormalTok{  \}}

\NormalTok{  p }\OtherTok{\textless{}{-}}\NormalTok{ jax\_scipy}\SpecialCharTok{$}\FunctionTok{expit}\NormalTok{(predictor)}
\SpecialCharTok{==}\ErrorTok{=====}
\FunctionTok{library}\NormalTok{(reticulate)}
\FunctionTok{library}\NormalTok{(BayesianInference)}

\NormalTok{m }\OtherTok{\textless{}{-}} \FunctionTok{importBI}\NormalTok{(}\AttributeTok{platform =} \StringTok{"cpu"}\NormalTok{)}
\NormalTok{jnp }\OtherTok{\textless{}{-}} \FunctionTok{import}\NormalTok{(}\StringTok{"jax.numpy"}\NormalTok{)}

\NormalTok{m}\SpecialCharTok{$}\NormalTok{data\_on\_model }\OtherTok{\textless{}{-}} \FunctionTok{list}\NormalTok{(}\AttributeTok{data =}\NormalTok{ data)}

\NormalTok{bi\_model\_duration }\OtherTok{\textless{}{-}} \ControlFlowTok{function}\NormalTok{(data) \{}

  \CommentTok{\# 1. Edge weights prior with partial pooling : Set the baseline (Intercept)}
\NormalTok{  edge\_sigma }\OtherTok{\textless{}{-}}\NormalTok{ m}\SpecialCharTok{$}\NormalTok{dist}\SpecialCharTok{$}\FunctionTok{half\_normal}\NormalTok{(}\FloatTok{2.0}\NormalTok{, }\AttributeTok{name =} \StringTok{"edge\_sigma"}\NormalTok{)}
\NormalTok{  edge\_weight }\OtherTok{\textless{}{-}}\NormalTok{ m}\SpecialCharTok{$}\NormalTok{dist}\SpecialCharTok{$}\FunctionTok{normal}\NormalTok{(}\FloatTok{0.0}\NormalTok{, edge\_sigma, }\AttributeTok{shape =} \FunctionTok{tuple}\NormalTok{(data}\SpecialCharTok{$}\NormalTok{num\_edges), }\AttributeTok{name =} \StringTok{"edge\_weight"}\NormalTok{)}
  
  \CommentTok{\# Additional absolute rate parameter strictly positive}
\NormalTok{  rate\_positive }\OtherTok{\textless{}{-}}\NormalTok{ m}\SpecialCharTok{$}\NormalTok{dist}\SpecialCharTok{$}\FunctionTok{half\_normal}\NormalTok{(}\FloatTok{1.0}\NormalTok{, }\AttributeTok{shape =} \FunctionTok{tuple}\NormalTok{(data}\SpecialCharTok{$}\NormalTok{num\_edges), }\AttributeTok{name =} \StringTok{"rate"}\NormalTok{)}
  
\NormalTok{  dyad\_ids\_0 }\OtherTok{\textless{}{-}}\NormalTok{ data}\SpecialCharTok{$}\NormalTok{dyad\_ids }\SpecialCharTok{{-}} \DecValTok{1}\NormalTok{L}
  \DocumentationTok{\#\# Build predictor}
\NormalTok{  predictor }\OtherTok{\textless{}{-}}\NormalTok{ edge\_weight[dyad\_ids\_0]}
  
  \CommentTok{\# 2. Fixed nodal factor(s)}
\NormalTok{  beta\_fixed }\OtherTok{\textless{}{-}}\NormalTok{ m}\SpecialCharTok{$}\NormalTok{dist}\SpecialCharTok{$}\FunctionTok{normal}\NormalTok{(}\FloatTok{0.0}\NormalTok{, }\FloatTok{1.0}\NormalTok{, }\AttributeTok{shape =} \FunctionTok{tuple}\NormalTok{(data}\SpecialCharTok{$}\NormalTok{num\_fixed), }\AttributeTok{name =} \StringTok{"beta\_fixed"}\NormalTok{)}
\NormalTok{  predictor }\OtherTok{\textless{}{-}}\NormalTok{ predictor }\SpecialCharTok{+}\NormalTok{ jnp}\SpecialCharTok{$}\FunctionTok{dot}\NormalTok{(data}\SpecialCharTok{$}\NormalTok{design\_fixed, beta\_fixed)}

  \CommentTok{\# 3. Random nodal effect(s)}
\NormalTok{  random\_group\_mu }\OtherTok{\textless{}{-}}\NormalTok{ m}\SpecialCharTok{$}\NormalTok{dist}\SpecialCharTok{$}\FunctionTok{normal}\NormalTok{(}\FloatTok{0.0}\NormalTok{, }\FloatTok{1.0}\NormalTok{, }\AttributeTok{shape =} \FunctionTok{tuple}\NormalTok{(data}\SpecialCharTok{$}\NormalTok{num\_random\_groups), }\AttributeTok{name =} \StringTok{"random\_group\_mu"}\NormalTok{)}
\NormalTok{  random\_group\_sigma }\OtherTok{\textless{}{-}}\NormalTok{ m}\SpecialCharTok{$}\NormalTok{dist}\SpecialCharTok{$}\FunctionTok{half\_normal}\NormalTok{(}\FloatTok{1.0}\NormalTok{, }\AttributeTok{shape =} \FunctionTok{tuple}\NormalTok{(data}\SpecialCharTok{$}\NormalTok{num\_random\_groups), }\AttributeTok{name =} \StringTok{"random\_group\_sigma"}\NormalTok{)}
  
\NormalTok{  group\_idx\_0 }\OtherTok{\textless{}{-}}\NormalTok{ data}\SpecialCharTok{$}\NormalTok{random\_group\_index }\SpecialCharTok{{-}} \DecValTok{1}\NormalTok{L}
\NormalTok{  beta\_random }\OtherTok{\textless{}{-}}\NormalTok{ m}\SpecialCharTok{$}\NormalTok{dist}\SpecialCharTok{$}\FunctionTok{normal}\NormalTok{(random\_group\_mu[group\_idx\_0], random\_group\_sigma[group\_idx\_0], }\AttributeTok{shape =} \FunctionTok{tuple}\NormalTok{(data}\SpecialCharTok{$}\NormalTok{num\_random), }\AttributeTok{name =} \StringTok{"beta\_random"}\NormalTok{)}
  
\NormalTok{  predictor }\OtherTok{\textless{}{-}}\NormalTok{ predictor }\SpecialCharTok{+}\NormalTok{ jnp}\SpecialCharTok{$}\FunctionTok{dot}\NormalTok{(data}\SpecialCharTok{$}\NormalTok{design\_random, beta\_random)}

  \CommentTok{\# 4. Link function}
\NormalTok{  p }\OtherTok{\textless{}{-}}\NormalTok{ m}\SpecialCharTok{$}\FunctionTok{link.inv\_logit}\NormalTok{(predictor)}
\SpecialCharTok{\textgreater{}}\ErrorTok{\textgreater{}\textgreater{}\textgreater{}\textgreater{}\textgreater{}\textgreater{}} \DecValTok{20}\NormalTok{c35a832b7f5f82f2011e8aad8fe0f1e2851c1b}
  
\NormalTok{  lambda\_exp }\OtherTok{\textless{}{-}}\NormalTok{ rate\_positive[dyad\_ids\_0] }\SpecialCharTok{/}\NormalTok{ p}
\NormalTok{  lambda\_pois }\OtherTok{\textless{}{-}}\NormalTok{ rate\_positive }\SpecialCharTok{*}\NormalTok{ data}\SpecialCharTok{$}\NormalTok{divisor}
  
\SpecialCharTok{\textless{}}\ErrorTok{\textless{}\textless{}\textless{}\textless{}\textless{}\textless{}}\NormalTok{ HEAD}
\NormalTok{  m}\SpecialCharTok{$}\NormalTok{dist}\SpecialCharTok{$}\FunctionTok{exponential}\NormalTok{(lambda\_exp, }\AttributeTok{obs=}\NormalTok{data}\SpecialCharTok{$}\NormalTok{event, }\AttributeTok{name=}\StringTok{"event"}\NormalTok{)}
\NormalTok{  m}\SpecialCharTok{$}\NormalTok{dist}\SpecialCharTok{$}\FunctionTok{poisson}\NormalTok{(lambda\_pois, }\AttributeTok{obs=}\NormalTok{data}\SpecialCharTok{$}\NormalTok{event\_count, }\AttributeTok{name=}\StringTok{"event\_count"}\NormalTok{)}
\SpecialCharTok{==}\ErrorTok{=====}
  \CommentTok{\# 5. Joint likelihood inference}
\NormalTok{  m}\SpecialCharTok{$}\NormalTok{dist}\SpecialCharTok{$}\FunctionTok{exponential}\NormalTok{(lambda\_exp, }\AttributeTok{obs =}\NormalTok{ data}\SpecialCharTok{$}\NormalTok{event, }\AttributeTok{name =} \StringTok{"event"}\NormalTok{)}
\NormalTok{  m}\SpecialCharTok{$}\NormalTok{dist}\SpecialCharTok{$}\FunctionTok{poisson}\NormalTok{(lambda\_pois, }\AttributeTok{obs =}\NormalTok{ data}\SpecialCharTok{$}\NormalTok{event\_count, }\AttributeTok{name =} \StringTok{"event\_count"}\NormalTok{)}
\SpecialCharTok{\textgreater{}}\ErrorTok{\textgreater{}\textgreater{}\textgreater{}\textgreater{}\textgreater{}\textgreater{}} \DecValTok{20}\NormalTok{c35a832b7f5f82f2011e8aad8fe0f1e2851c1b}
\NormalTok{\}}

\NormalTok{m}\SpecialCharTok{$}\FunctionTok{fit}\NormalTok{(bi\_model\_duration)}
\NormalTok{m}\SpecialCharTok{$}\FunctionTok{summary}\NormalTok{()}
\end{Highlighting}
\end{Shaded}

\subsection{Julia}

\begin{Shaded}
\begin{Highlighting}[]
\ImportTok{using} \BuiltInTok{BayesianInference}

\NormalTok{m }\OperatorTok{=} \FunctionTok{importBI}\NormalTok{(platform}\OperatorTok{=}\StringTok{"cpu"}\NormalTok{)}
\OperatorTok{\textless{}\textless{}\textless{}\textless{}\textless{}\textless{}\textless{}}\NormalTok{ HEAD}
\NormalTok{m.data\_on\_model }\OperatorTok{=} \FunctionTok{Dict}\NormalTok{(}\StringTok{"data"} \OperatorTok{=\textgreater{}}\NormalTok{ data\_list\_jax)}

\PreprocessorTok{@BI} \KeywordTok{function} \FunctionTok{bi\_model\_duration}\NormalTok{(data)}
\NormalTok{    edge\_weight }\OperatorTok{=}\NormalTok{ m.dist.}\FunctionTok{normal}\NormalTok{(data[}\StringTok{"prior\_edge\_mu"}\NormalTok{], data[}\StringTok{"prior\_edge\_sigma"}\NormalTok{], }
\NormalTok{                                shape}\OperatorTok{=}\NormalTok{(data[}\StringTok{"num\_edges"}\NormalTok{],), name}\OperatorTok{=}\StringTok{"edge\_weight"}\NormalTok{)}
    
\NormalTok{    rate\_param }\OperatorTok{=}\NormalTok{ m.dist.}\FunctionTok{normal}\NormalTok{(}\FloatTok{0.0}\NormalTok{, data[}\StringTok{"prior\_rate\_sigma"}\NormalTok{], }
\NormalTok{                               shape}\OperatorTok{=}\NormalTok{(data[}\StringTok{"num\_edges"}\NormalTok{],), name}\OperatorTok{=}\StringTok{"rate"}\NormalTok{)}
\NormalTok{    rate\_positive }\OperatorTok{=}\NormalTok{ jnp.}\FunctionTok{abs}\NormalTok{(rate\_param)}
    
\NormalTok{    dyad\_ids\_0 }\OperatorTok{=}\NormalTok{ data[}\StringTok{"dyad\_ids"}\NormalTok{] }\OperatorTok{.{-}} \FloatTok{1}
\NormalTok{    predictor }\OperatorTok{=}\NormalTok{ edge\_weight[dyad\_ids\_0]}
    
    \ControlFlowTok{if}\NormalTok{ data[}\StringTok{"num\_fixed"}\NormalTok{] }\OperatorTok{\textgreater{}} \FloatTok{0}
\NormalTok{        beta\_fixed }\OperatorTok{=}\NormalTok{ m.dist.}\FunctionTok{normal}\NormalTok{(data[}\StringTok{"prior\_fixed\_mu"}\NormalTok{], data[}\StringTok{"prior\_fixed\_sigma"}\NormalTok{], }
\NormalTok{                                   shape}\OperatorTok{=}\NormalTok{(data[}\StringTok{"num\_fixed"}\NormalTok{],), name}\OperatorTok{=}\StringTok{"beta\_fixed"}\NormalTok{)}
\NormalTok{        predictor }\OperatorTok{=}\NormalTok{ predictor }\OperatorTok{.+}\NormalTok{ jnp.}\FunctionTok{dot}\NormalTok{(data[}\StringTok{"design\_fixed"}\NormalTok{], beta\_fixed)}
    \ControlFlowTok{end}

\NormalTok{    p }\OperatorTok{=}\NormalTok{ jax.scipy.special.}\FunctionTok{expit}\NormalTok{(predictor)}
\OperatorTok{=======}
\NormalTok{jnp }\OperatorTok{=}\NormalTok{ m.jax.numpy}

\NormalTok{m.data\_on\_model }\OperatorTok{=} \FunctionTok{Dict}\NormalTok{(}\StringTok{"data"} \OperatorTok{=\textgreater{}}\NormalTok{ data)}

\PreprocessorTok{@BI} \KeywordTok{function} \FunctionTok{bi\_model\_duration}\NormalTok{(data)}

    \CommentTok{\# 1. Edge weights prior with partial pooling : Set the baseline (Intercept)}
\NormalTok{    edge\_sigma }\OperatorTok{=}\NormalTok{ m.dist.}\FunctionTok{half\_normal}\NormalTok{(}\FloatTok{2.0}\NormalTok{, name}\OperatorTok{=}\StringTok{"edge\_sigma"}\NormalTok{)}
\NormalTok{    edge\_weight }\OperatorTok{=}\NormalTok{ m.dist.}\FunctionTok{normal}\NormalTok{(}\FloatTok{0.0}\NormalTok{, edge\_sigma, shape}\OperatorTok{=}\NormalTok{(data[}\StringTok{"num\_edges"}\NormalTok{],), name}\OperatorTok{=}\StringTok{"edge\_weight"}\NormalTok{)}
    
    \CommentTok{\# Additional absolute rate parameter strictly positive}
\NormalTok{    rate\_positive }\OperatorTok{=}\NormalTok{ m.dist.}\FunctionTok{half\_normal}\NormalTok{(}\FloatTok{1.0}\NormalTok{, shape}\OperatorTok{=}\NormalTok{(data[}\StringTok{"num\_edges"}\NormalTok{],), name}\OperatorTok{=}\StringTok{"rate"}\NormalTok{)}
    
\NormalTok{    dyad\_ids\_0 }\OperatorTok{=}\NormalTok{ data[}\StringTok{"dyad\_ids"}\NormalTok{] }\OperatorTok{.{-}} \FloatTok{1}
    \CommentTok{\#\# Build predictor}
\NormalTok{    predictor }\OperatorTok{=}\NormalTok{ edge\_weight[dyad\_ids\_0]}
    
    \CommentTok{\# 2. Fixed nodal factor(s)}
\NormalTok{    beta\_fixed }\OperatorTok{=}\NormalTok{ m.dist.}\FunctionTok{normal}\NormalTok{(}\FloatTok{0.0}\NormalTok{, }\FloatTok{1.0}\NormalTok{, shape}\OperatorTok{=}\NormalTok{(data[}\StringTok{"num\_fixed"}\NormalTok{],), name}\OperatorTok{=}\StringTok{"beta\_fixed"}\NormalTok{)}
\NormalTok{    predictor }\OperatorTok{=}\NormalTok{ predictor }\OperatorTok{.+}\NormalTok{ jnp.}\FunctionTok{dot}\NormalTok{(data[}\StringTok{"design\_fixed"}\NormalTok{], beta\_fixed)}

    \CommentTok{\# 3. Random nodal effect(s)}
\NormalTok{    random\_group\_mu }\OperatorTok{=}\NormalTok{ m.dist.}\FunctionTok{normal}\NormalTok{(}\FloatTok{0.0}\NormalTok{, }\FloatTok{1.0}\NormalTok{, shape}\OperatorTok{=}\NormalTok{(data[}\StringTok{"num\_random\_groups"}\NormalTok{],), name}\OperatorTok{=}\StringTok{"random\_group\_mu"}\NormalTok{)}
\NormalTok{    random\_group\_sigma }\OperatorTok{=}\NormalTok{ m.dist.}\FunctionTok{half\_normal}\NormalTok{(}\FloatTok{1.0}\NormalTok{, shape}\OperatorTok{=}\NormalTok{(data[}\StringTok{"num\_random\_groups"}\NormalTok{],), name}\OperatorTok{=}\StringTok{"random\_group\_sigma"}\NormalTok{)}
    
\NormalTok{    group\_idx\_0 }\OperatorTok{=}\NormalTok{ data[}\StringTok{"random\_group\_index"}\NormalTok{] }\OperatorTok{.{-}} \FloatTok{1}
\NormalTok{    beta\_random }\OperatorTok{=}\NormalTok{ m.dist.}\FunctionTok{normal}\NormalTok{(random\_group\_mu[group\_idx\_0], random\_group\_sigma[group\_idx\_0], shape}\OperatorTok{=}\NormalTok{(data[}\StringTok{"num\_random"}\NormalTok{],), name}\OperatorTok{=}\StringTok{"beta\_random"}\NormalTok{)}
    
\NormalTok{    predictor }\OperatorTok{=}\NormalTok{ predictor }\OperatorTok{.+}\NormalTok{ jnp.}\FunctionTok{dot}\NormalTok{(data[}\StringTok{"design\_random"}\NormalTok{], beta\_random)}

    \CommentTok{\# 4. Link function}
\NormalTok{    p }\OperatorTok{=}\NormalTok{ m.link.}\FunctionTok{inv\_logit}\NormalTok{(predictor)}
\OperatorTok{\textgreater{}\textgreater{}\textgreater{}\textgreater{}\textgreater{}\textgreater{}\textgreater{}} \FloatTok{20}\NormalTok{c35a832b7f5f82f2011e8aad8fe0f1e2851c1b}
    
\NormalTok{    lambda\_exp }\OperatorTok{=}\NormalTok{ rate\_positive[dyad\_ids\_0] }\OperatorTok{./}\NormalTok{ p}
\NormalTok{    lambda\_pois }\OperatorTok{=}\NormalTok{ rate\_positive }\OperatorTok{.*}\NormalTok{ data[}\StringTok{"divisor"}\NormalTok{]}
    
\OperatorTok{\textless{}\textless{}\textless{}\textless{}\textless{}\textless{}\textless{}}\NormalTok{ HEAD}
\OperatorTok{=======}
    \CommentTok{\# 5. Joint likelihood inference}
\OperatorTok{\textgreater{}\textgreater{}\textgreater{}\textgreater{}\textgreater{}\textgreater{}\textgreater{}} \FloatTok{20}\NormalTok{c35a832b7f5f82f2011e8aad8fe0f1e2851c1b}
\NormalTok{    m.dist.}\FunctionTok{exponential}\NormalTok{(lambda\_exp, obs}\OperatorTok{=}\NormalTok{data[}\StringTok{"event"}\NormalTok{], name}\OperatorTok{=}\StringTok{"event"}\NormalTok{)}
\NormalTok{    m.dist.}\FunctionTok{poisson}\NormalTok{(lambda\_pois, obs}\OperatorTok{=}\NormalTok{data[}\StringTok{"event\_count"}\NormalTok{], name}\OperatorTok{=}\StringTok{"event\_count"}\NormalTok{)}
\KeywordTok{end}

\NormalTok{m.}\FunctionTok{fit}\NormalTok{(bi\_model\_duration)}
\NormalTok{m.}\FunctionTok{summary}\NormalTok{()}
\end{Highlighting}
\end{Shaded}

\textless\textless\textless\textless\textless\textless\textless{} HEAD
\#\#\#\# Mathematical Details ======= \#\#\# Mathematical Details
\textgreater\textgreater\textgreater\textgreater\textgreater\textgreater\textgreater{}
20c35a832b7f5f82f2011e8aad8fe0f1e2851c1b The mean behavior time
\(\mu_{ij}\) can be recovered from the estimated proportion scale models
representing expected duration lengths per sequence generated by
overarching absolute rates \(\lambda_{ij}\). \[
X_{ij}^{(n)} \sim \text{Exponential}(\lambda_{ij} / t_{ij}^{(n)}) 
\] \[
K_{ij} \sim \text{Poisson}(\lambda_{ij} D_{ij})
\] The proportion is generated matching logistic mapping inputs: \[
\text{logit}(t_{ij}^{(n)}) = \omega_{ij} + \dots 
\] And recovered as expected proportion values.

\paragraph{Notes}\label{notes-2}

\begin{tcolorbox}[enhanced jigsaw, colframe=quarto-callout-note-color-frame, bottomtitle=1mm, coltitle=black, titlerule=0mm, colbacktitle=quarto-callout-note-color!10!white, colback=white, opacitybacktitle=0.6, rightrule=.15mm, toptitle=1mm, opacityback=0, arc=.35mm, left=2mm, title=\textcolor{quarto-callout-note-color}{\faInfo}\hspace{0.5em}{Note}, breakable, bottomrule=.15mm, toprule=.15mm, leftrule=.75mm]

Because edge estimates define probability scales while absolute
intensity generates sequence frequency, parameter isolation allows
complex uncoupled insights such as distinct testing comparing the
absolute magnitude of interactions versus simply calculating engaged
durations.

\end{tcolorbox}

\section{\textless\textless\textless\textless\textless\textless\textless{}
HEAD}\label{head-1}

\subsubsection{Notes}\label{notes-3}

\begin{tcolorbox}[enhanced jigsaw, colframe=quarto-callout-note-color-frame, bottomtitle=1mm, coltitle=black, titlerule=0mm, colbacktitle=quarto-callout-note-color!10!white, colback=white, opacitybacktitle=0.6, rightrule=.15mm, toptitle=1mm, opacityback=0, arc=.35mm, left=2mm, title=\textcolor{quarto-callout-note-color}{\faInfo}\hspace{0.5em}{Note}, breakable, bottomrule=.15mm, toprule=.15mm, leftrule=.75mm]

Varying effects for nodes ad dyads could be build using Multi Variate
Normal distribution to account for Node-Level and/or dyadic Correlated
Effects.

\begin{Shaded}
\begin{Highlighting}[]
  \ControlFlowTok{if}\NormalTok{ (}\ImportTok{as}\NormalTok{.numeric(data$num\_random) }\OperatorTok{\textgreater{}} \DecValTok{0}\NormalTok{) \{}
    \CommentTok{\# Extract group ID from design\_random matrix}
    \CommentTok{\# Assumes design\_random is a one{-}hot matrix (N\_obs, N\_groups)}
\NormalTok{    group\_ids\_0 }\OperatorTok{\textless{}{-}}\NormalTok{ jnp$argmax(data$design\_random, axis }\OperatorTok{=} \DecValTok{1L}\NormalTok{)}

    \ControlFlowTok{if}\NormalTok{ (}\ImportTok{as}\NormalTok{.numeric(data$num\_random) }\OperatorTok{==} \DecValTok{1L}\NormalTok{) \{}
\NormalTok{      varying\_intercepts }\OperatorTok{\textless{}{-}}\NormalTok{ m$effects$varying\_intercept(}
\NormalTok{        N\_groups }\OperatorTok{=} \ImportTok{as}\NormalTok{.integer(data$num\_random\_groups),}
\NormalTok{        group\_id }\OperatorTok{=}\NormalTok{ group\_ids\_0,}
\NormalTok{        a\_bar }\OperatorTok{=}\NormalTok{ m$dist$normal(}\FloatTok{0.0}\NormalTok{, }\FloatTok{1.0}\NormalTok{, shape }\OperatorTok{=} \BuiltInTok{tuple}\NormalTok{(}\DecValTok{1L}\NormalTok{), name }\OperatorTok{=} \StringTok{"random\_group\_mu\_new"}\NormalTok{),}
\NormalTok{        sigma }\OperatorTok{=}\NormalTok{ m$dist$exponential(}\FloatTok{1.0}\NormalTok{, shape }\OperatorTok{=} \BuiltInTok{tuple}\NormalTok{(}\DecValTok{1L}\NormalTok{), name }\OperatorTok{=} \StringTok{"random\_group\_sigma\_new"}\NormalTok{),}
\NormalTok{        group\_name }\OperatorTok{=} \StringTok{"node\_id"}
\NormalTok{      )}
\NormalTok{      predictor }\OperatorTok{\textless{}{-}}\NormalTok{ predictor }\OperatorTok{+}\NormalTok{ varying\_intercepts}
\NormalTok{    \} }\ControlFlowTok{else}\NormalTok{ \{}
      \CommentTok{\# Use varying\_effects if we have more than 1 effect (intercept + slopes)}
\NormalTok{      out }\OperatorTok{\textless{}{-}}\NormalTok{ m$effects$varying\_effects(}
\NormalTok{        N\_vars }\OperatorTok{=} \ImportTok{as}\NormalTok{.integer(data$num\_random) }\OperatorTok{{-}} \DecValTok{1L}\NormalTok{,}
\NormalTok{        N\_group }\OperatorTok{=} \ImportTok{as}\NormalTok{.integer(data$num\_random\_groups),}
\NormalTok{        group\_id }\OperatorTok{=}\NormalTok{ group\_ids\_0,}
\NormalTok{        group\_name }\OperatorTok{=} \StringTok{"node\_id"}\NormalTok{,}
\NormalTok{        alpha\_bar }\OperatorTok{=}\NormalTok{ m$dist$normal(}\FloatTok{0.0}\NormalTok{, }\FloatTok{1.0}\NormalTok{, shape }\OperatorTok{=} \BuiltInTok{tuple}\NormalTok{(}\DecValTok{1L}\NormalTok{), name }\OperatorTok{=} \StringTok{"random\_group\_mu\_new"}\NormalTok{),}
\NormalTok{        sd\_intercept }\OperatorTok{=}\NormalTok{ m$dist$exponential(}\FloatTok{1.0}\NormalTok{, shape }\OperatorTok{=} \BuiltInTok{tuple}\NormalTok{(}\DecValTok{1L}\NormalTok{), name }\OperatorTok{=} \StringTok{"random\_group\_sigma\_new"}\NormalTok{)}
\NormalTok{      )}
\NormalTok{      predictor }\OperatorTok{\textless{}{-}}\NormalTok{ predictor }\OperatorTok{+}\NormalTok{ out[[}\DecValTok{1}\NormalTok{]]}
      \CommentTok{\# If slopes are present, they are in out[[2]], but bisonR standard model used here}
      \CommentTok{\# usually doesn\textquotesingle{}t have slopes in design\_random yet in this simple way}
\NormalTok{    \}}
\NormalTok{  \}}
\end{Highlighting}
\end{Shaded}

\end{tcolorbox}

\begin{quote}
\begin{quote}
\begin{quote}
\begin{quote}
\begin{quote}
\begin{quote}
\begin{quote}
20c35a832b7f5f82f2011e8aad8fe0f1e2851c1b
\end{quote}
\end{quote}
\end{quote}
\end{quote}
\end{quote}
\end{quote}
\end{quote}

\phantomsection\label{refs}
\begin{CSLReferences}{1}{0}
\bibitem[\citeproctext]{ref-BISoN}
Hart, Jordan, Michael Nash Weiss, Daniel Franks, and Lauren Brent. 2023.
{``BISoN: A Bayesian Framework for Inference of Social Networks.''}
\emph{Methods in Ecology and Evolution} 14: 2411--20.
https://doi.org/\url{https://doi.org/10.1111/2041-210X.14171}.

\end{CSLReferences}




\end{document}
